%Chargement des paquets
\usepackage{amsmath}
\usepackage{amsfonts}
\usepackage{amssymb}
\usepackage{enumerate}
\usepackage{mathtools}
\usepackage{bbm}
\usepackage{xparse, etoolbox}
\usepackage{enumerate}
\usepackage{enumitem}
\usepackage{mathabx}
\usepackage{babel}[french]

%environment
\usepackage{amsthm}
\usepackage{keytheorems}
\usepackage{hyperref} 

%Interface théorème
\newkeytheorem{lem}[
	name = Lemme
]

\newkeytheorem{prop}[
	name = Proposition
]

\newkeytheorem{thm}[
	name = Théorème
]

\newkeytheorem{coro}[
	name = Corollaire
]

\renewcommand*{\proofname}{Démonstration}

\theoremstyle{definition}
\newtheorem{defn}{Définition}

\theoremstyle{definition}
\newtheorem*{nota}{Notation}

\theoremstyle{definition}
\newtheorem*{limi}{Liminaire}

\theoremstyle{remark}
\newtheorem{rmq}{Remarque}

\theoremstyle{plain}
\newtheorem*{fait}{Fait}


%Ensembles de nombres
\newcommand{\N} {\mathbb{N}}
\newcommand{\Ne}{\N^\ast}
\newcommand{\Z} {\mathbb{Z}}
\newcommand{\Q} {\mathbb{Q}}
\newcommand{\R} {\mathbb{R}}
\newcommand{\Rb}{\overline{\mathbb{R}}}
\newcommand{\Rp}{\R_+}
\newcommand{\Rm}{\R_-}
\newcommand{\K} {\mathbb{K}}
\newcommand{\Ket} {\mathbb{K^{\ast}}}
\newcommand{\Cx}{\mathbb{C}}

%Ensembles, tribus
\newcommand{\A}{\mathbb{A}}
\renewcommand{\P}{\mathcal{P}}
\newcommand{\T}{\mathcal{T}}
\newcommand{\F}{\mathcal{F}}
\newcommand{\C} {\mathcal{C}}
\newcommand{\ouv}{\mathcal{O}}
\newcommand{\B}{\mathcal{B}}
\newcommand{\M}{\mathcal{M}}
\newcommand{\etag}{\mathcal{E}}
\newcommand{\etagOp}{\etag^{+}(\Omega)}
\newcommand{\etagO}{\etag(\Omega)}
%%Lp avant quotientage
\newcommand{\Lic}{\mathcal{L}^1}
\newcommand{\Lpc}{\mathcal{L}^p}
\newcommand{\Linfc}{\mathcal{L}^{\infty}}
\newcommand{\Lifull}{\Lic(\Omega, \B, m)}
\newcommand{\Lpfull}{\Lpc(\Omega, \B, m)}
\newcommand{\Linffull}{\Linfc(\Omega, \B, m)}
%%Lp après quotientage
\newcommand{\Li}{L^1}
\newcommand{\Ld}{L^2}
\newcommand{\Lp}{L^p}
\newcommand{\Linf}{L^{\infty}}
\newcommand{\Listd}{\Li(\Omega, m)} 
\newcommand{\Ldstd}{\Ld(\Omega, m)} 
\newcommand{\Lpstd}{\Lp(\Omega, m)} 

%Quelques opérateurs
\DeclareMathOperator{\codim}{codim}
\DeclareMathOperator{\rg}{rg}
\DeclareMathOperator{\tr}{tr}
\DeclareMathOperator{\vect}{Vect}
\DeclareMathOperator{\ima}{Im}
\DeclareMathOperator{\id}{Id}
\DeclareMathOperator{\sign}{sign}
\DeclareMathOperator{\ext}{ext}
\newcommand{\ind}{\mathbbm{1}}
\newcommand*{\spr}[2]{\left(\,#1\,\middle|\,#2\,\right)}
\newcommand{\interior}[1]{%
  {\kern0pt#1}^{\mathrm{o}}%
}
\DeclareMathOperator{\card}{card}
\DeclareMathOperator{\ppcm}{ppcm}

%Commandes de pure flemme
\newcommand{\veps}{\varepsilon}
\newcommand{\intab}{\int_{a}^{b}}
\newcommand{\intR}{\int_{-\infty}^{\infty}}
\newcommand{\intinf}{\int_{- \infty}^{+ \infty}}
\newcommand{\equi}{\Leftrightarrow}
\newcommand{\Kev}{$\K$-espace vectoriel }
\newcommand{\Kevs}{$\K$-espaces vectoriels }
\newcommand{\Rev}{$\R$-espace vectoriel }
\newcommand{\Revs}{$\R$-espaces vectoriels }
\newcommand{\Cev}{$\Cx$-espace vectoriel }
\newcommand{\Cevs}{$\Cx$-espaces vectoriels }
\newcommand{\evn}{(E, \norm{.})}
\newcommand{\interf}[2]{[#1,#2]}
\newcommand{\limb}{\lim\limits_}
\newcommand{\lebn}{\lambda_n}
\newcommand{\pp}{\text{~presque partout}}
\newcommand{\derivp}[2]{\frac{\partial #1}{\partial #2}}
\newcommand{\famiu}{(u_i)_{i \in I}}
\newcommand{\sev}{\text{~s.e.v.}}
\newcommand{\Mnp}{M_{n,p}(\K)}
\newcommand{\Mn}{M_{n}(\K)}
\newcommand{\tq}{\text{~t.q.~}}
\newcommand{\mq}{~montrer~que~}
\newcommand{\et}{\quad \text{et} \quad}
\newcommand{\ou}{\text{~ou~}}
\newcommand{\Kx}{\K[X]}


%Commandes de gars chelou
\renewcommand{\subset}{\subseteq}

%Commande restriction, ty duckduckgo
\def\restr#1#2{\mathchoice
              {\setbox1\hbox{${\displaystyle #1}_{\scriptstyle #2}$}
              \restraux{#1}{#2}}
              {\setbox1\hbox{${\textstyle #1}_{\scriptstyle #2}$}
              \restraux{#1}{#2}}
              {\setbox1\hbox{${\scriptstyle #1}_{\scriptscriptstyle #2}$}
              \restraux{#1}{#2}}
              {\setbox1\hbox{${\scriptscriptstyle #1}_{\scriptscriptstyle #2}$}
              \restraux{#1}{#2}}}
\def\restraux#1#2{{#1\,\smash{\vrule height .8\ht1 depth .85\dp1}}_{\,#2}} 

%Quelques normes
\DeclarePairedDelimiter\abs{\lvert}{\rvert}
\DeclarePairedDelimiter\labs{\bigl\lvert}{\bigr\rvert}
\DeclarePairedDelimiter\norm{\lVert}{\rVert}
\DeclarePairedDelimiterXPP\normi[1]{}\lVert\rVert{_1}{\ifblank{#1}{\:\cdot\:}{#1}}
\DeclarePairedDelimiterXPP\normd[1]{}\lVert\rVert{_2}{\ifblank{#1}{\:\cdot\:}{#1}}
\DeclarePairedDelimiterXPP\normp[1]{}\lVert\rVert{_p}{\ifblank{#1}{\:\cdot\:}{#1}}
\DeclarePairedDelimiterXPP\normq[1]{}\lVert\rVert{_q}{\ifblank{#1}{\:\cdot\:}{#1}}
\DeclarePairedDelimiterXPP\norminf[1]{}\lVert\rVert{_\infty}{\ifblank{#1}{\:\cdot\:}{#1}}

%Bracket en angle
\DeclarePairedDelimiter\angl{\langle}{\rangle}

%Set, parenthèses
\newcommand{\set}[1]{\{ #1 \}}
\newcommand{\bigset}[1]{\bigl\{ #1 \bigr\}}
\newcommand{\Bigset}[1]{\Bigl\{ #1 \Bigr\}}
\newcommand{\bigpar}[1]{\bigl( #1 \bigr)}
\newcommand{\Bigpar}[1]{\Bigl( #1 \Bigr)}

%Opitmisation des opérateurs de comparaison
\DeclarePairedDelimiter\tio{o (}{)}
\DeclarePairedDelimiter\gro{O (}{)}


\pagestyle{fancy}

\fancyhead[C]{}
\fancyhead[R]{}

\begin{document}

\underline{Source :} 

\begin{itemize}
\item Tauvel - Analyse complexe - page 52
\item Hassan - Topologie - page 169 (pour la connexité)
\end{itemize}

\underline{Leçons :}
\begin{itemize}
\item
\end{itemize}

\begin{defn}
Un espace topologique $X$ est convexe s'il n'est pas réunion de deux ensembles ouverts non vides disjoints.
\end{defn}

\begin{thm}[Caractérisation de la connexité]
Soit $X$ un espace topologique, les propriétés suivants sont équivalentes :
\begin{enumerate}
\item $X$ est connexe ;
\item $X$ n'est pas réunion de deux fermés non vides disjoints ;
\item il n'existe pas dans $X$ d'autres parties ouvertes et fermées que $X$ et $\emptyset$ ;
\item toute application continue de $X$ dans l'espace discret $\Z$ est constante ;
\item toute application continue de $X$ dans l'espace discret $\set{0,1}$ est constante.
\end{enumerate}
\end{enumerate}
\end{thm}

\begin{proof}
$(i) \Rightarrow (ii)$ Supposons que $X = F \cup G$, avec $F, G$ fermés non vides disjoints.
Clairement, $(X \setminus F) \cup (X \setminus G) = F$ et les ensembles du membre de gauche sont des ouverts, ce qui est contradictoire. \\

$(ii) \Rightarrow (iii)$ Soit $\emptyset \subsetneq F \subsetneq X$, avec $F$ ouvert et fermé. 
Clairement $X= F \cup (X \setminus F)$ réunion de deux fermés. \\

$(iii) \Rightarrow (iv)$ Soit $x_0 \in X$ et $f(x_0) \in \Z$. L'ensemble $\set{f(x_0)}$ est ouvert et fermé dans $\Z$, 
donc son image réciproque par $f$ est un ouvert et fermé non vide de $X$, c'est donc $X$, $f$ est constante.

$(iv) \Rightarrow (v)$ On considère $i$ l'injection canonique de $\set{0,1}$ dans $\Z$. Soit $f$ continue de $X$ dans $\set{0,1}$,
alors $i \circ f$ est continue donc constante, ce qui implique que $f$ est constante.

$(v) \Rightarrow (i)$ Supposons $X$ non connexe, soit $X = U \cup V$ avec $U, V$ des ouverts non vides de $X$.
L'application $f$ définie sur $X$ par $f(x) = 0$ si $x \in U$ ; $f(x) = 1$ sinon, est continue donc constante, ce qui est absurde.
\end{proof}

\begin{thm}[Principe du prolongement analytique]
\label{thm:1}
Soient $U$ un ouvert connexe de $\Cx$, a \in U$ et $f$ une fonction analytique sur $U$. Les conditions suivantes sont équivalentes :
\begin{enumerate}[label=(\roman*)]
\item $f$ est identiquement nulle sur $U$ ;
\item $f$ est identiquement nulle sur un voisinage de $a$ ;
\item pour tout $n \in \N$, on a $f^{(n)}(a) = 0$.
\end{enumerate}
\end{thm}

\begin{proof}
Les implications $(i) \Rightarrow (ii) \Rightarrow (iii)$ sont claires. \\

$(iii) \Rightarrow (ii)$ Par définition, il existe un voisinage de $a$ sur lequel $f$ est égale à son développement en série entière,
donc identiquement nulle sur ce voisinage. \\

$(ii) \Rightarrow (i)$ Soit $V$ l'ensemble des $z \in U$ tel que $f$ soit identiquement nulle dans un voisinage de $z$.
$V$ est non vide par hypothèse et si $z \in V$, alors il existe un ouvert contenant $z$ et sur lequel $f$ est identiquement nulle.
Donc cet ouvert est dans $V$ ce qui démontre que $V$ est ouvert.
Prenons $(z_n)_\N$ une suite de points de $V$ convergeant vers $z \in U$. 
Pour tous $n,k$ \in \N$, on a $f^{k}(z_n)=0$. Par continuité de chaque $f^{k}$, il vient que $f^{k}(z)=0$ pour tout $k \in \N$.
Comme $f$ est développable en série entière autour de $z$, on en déduit que $f$ est identiquement nulle au voisinage de $z$, soit 
$z \in V$ et $V$ est fermé. Par connexité de $U$, on déduit que $U=V$ et la fonction est identiquement nulle sur $U$.
\end{proof}

\begin{coro}
Soient $U$ un ouvert connexe de $\Cx$ e $f, g$ analytiques sur $U$. 
Si $f$ et $g$ coïncident au voisinage d'un point de $U$ alors $f$ et $g$ sont égales sur tout $U$.
\end{coro}

\begin{defn}[partie discrète]
Soient $U$ un ouvert de $\Cx$ et $A$ une partie de $U$. On dit que $A$ est une partie discrète de $U$ s, pour tout $a \in U$, 
il existe $r>0$ tel que $D(a,r) \cap A = \set{a}$.
\end{defn}

\begin{thm}
\label{thm:2}
Soient $U$ un ouvert de $\Cx$ et $A$ une partie de $U$. Les propositions suivantes sont équivalentes.
\begin{enumerate}[label=(\roman*)]
\item Tout $z \in U$ possède un voisinage $V$ tel que $V \cap A$ soit fini ;
\item pour tout compact $K$ de $U$, l'ensemble $K \cap A$ est fini ;
\iem $A$ est une partie discrète et fermée de $U$.
\end{enumerate}
Si ces conditions sont vérifiées, on dit que $A$ est une partie localement finie de $U$.
\end{thm}

\begin{proof}
$(i) \Rightarrow (ii)$ Pour $z \in K$ on note $V_z$ un voisinage de $z$ tel que $V \cap A$ soit fini.
Alors $K \subset \bigcup_{z} V_z$ et on extrait de de recouvrement un sous recouvrement fini. 
Ainsi, $K \cap A$ est fini. \\

$(ii) \Rightarrow (iii)$ Comme $\Cx$ est localement compact, tout point $z \in U$ possède un voisinage compact que l'on note $V_z$,
et $V_z \cap A$ est fini. 
Par séparation de $\Cx$, chaque point de $V_z \cap A$ admet un voisinage disjoint des voisinages des autres points, on peut donc
exhiber un voisinage $W_z$ de $z$ tel que $W_z \cap A = \set{z}$ donc $A$ est une partie discrète. \\
Prenons maintenant $z \in \overline{A}$, l'adhérence de $A$ dans $U$, si $z \not \in A$, alors le voisinage $W_z$ précédemment
exhibé est tel que $W_z \cap A = \emptyset$ ce qui est contradictoire. Donc $A$ est fermé. \\

$(iii) \Rightarrow (i)$ Prenons $z \in U$, si $z \in A$ alors il existe un voisinage $V$ de $z$ tel que $V \cap A = \se{z}$ car $A$
est une partie discrète. Sinon, comme $A$ est fermé, il existe un voisinage $V$ de $z$ tel que $V \cap A = \emptyset$.
\end{proof}

\begin{thm}[Principe des zéros isolés]
Soient $U$ un ouvert connexe de $\Cx$ et $f$ une fonction analytique sur $U$ non identiquement nulle. 
L'ensemble $Z(f)$ des zéros de $f$ est une partie localement finie de $U$.
\end{thm}

\begin{proof}
$Z(f) = f^{-1}(\set{0})$, c'est donc un fermé par continuité de $f$. 
Il reste donc à montrer que $Z(f)$ est une partie discrète de $U$ (cf. théorème $\ref{thm:2}$). 
Soit $u \in U$, d'après le théorème $\ref{thm:1}$, il existe $k \in \N$ tel que $f^{k}(u) \neq 0$. 
On note $k$ le plus petit de ces entiers, au voisinage de $u$, on a :
\[ f(z) = a_k (z-u)^k + (z-u)^k \sum_{n=1}^{\infty} a_{n+k} (z-u)^n , \]
avec $a_k \neq 0$. On définit la fonction $g$ sur ce voisinage par :
\[ g(z) = \sum_{n=1}^{\infty} a_{n+k} (z-u)^n . \]
Cette fonction est nulle en $u$ et développable en série entière en ce point, et en particulier continue en $u$.
Donc, il existe un voisinage $V$ de $u$ sur lequel $\abs{g(z)} \leq \abs{a_k}$.
Ainsi, pour tout $z \in V \setminus \set{z}$, on a :
\[ \abs{f(z)} \ leq \abs{z-u}^k ( \abs{a_k} - \abs{g(z)} > 0 . \]
Ainsi, $Z(f) \cap V = \set{u}$, ce qui achève la démonstration.
\end{proof}



\end{document}